\documentclass[UTF8,a4paper,11pt,openany]{ctexbook}
\usepackage{../common/statlearnbook}
\SLSetBookMetadata{第 4 册}{统计学习方法讲义 v2.6.0}{无监督学习与矩阵分解}
\begin{document}
\setcounter{page}{672}
\setcounter{chapter}{29}
\setcounter{figure}{1}
\pagestyle{headings}
\chaptermark{概率潜在语义分析}
\sectionmark{EM 参数估计}
\begin{geometrybox}
\textbf{几何解释。}每个$\phi_{:k}$是$(M-1)$维单词单纯形中的主题点。文档分布$P(w\mid d_j)=\Phi\theta_{:j}$是这些主题点的凸组合，因此全部训练文档都落在由主题点张成的凸包中。
\end{geometrybox}
% v2.6.0 figure UID: FIG-P525-01
\tikzset{
  slfig-FIG-P525-01/.style={font=\fontsize{9.4pt}{11.2pt}\selectfont,
    every path/.append style={line cap=round,line join=round}},
  slfig-FIG-P525-01-outer/.style={draw=SLTextGray,line width=.92pt},
  slfig-FIG-P525-01-hull/.style={draw=SLTeal,line width=.86pt},
  slfig-FIG-P525-01-spoke/.style={draw=SLGray,densely dashed,line width=.54pt},
  slfig-FIG-P525-01-weight/.style={fill=none,inner sep=0pt},
  slfig-FIG-P525-01-doclabel/.style={anchor=west,text=SLGold,fill=white,
    fill opacity=.92,text opacity=1,inner xsep=1pt,inner ysep=.5pt},
  slfig-FIG-P525-01-formula/.style={draw=SLRuleGray,rounded corners=2pt,
    fill=SLSoftGray,line width=.62pt,align=center,inner ysep=2.2pt},
  slfig-FIG-P525-01-legend/.style={anchor=west,text=SLTextGray,
    font=\fontsize{8.8pt}{10.4pt}\selectfont}
}
\begin{figure}[H]
\centering
\begin{tikzpicture}[slfig-FIG-P525-01,x=1cm,y=1cm]
  % 显式重心变换：(p_1,p_2,p_3) ->
  % (x,y)=(6(p_2+p_3/2), 6(sqrt(3)/2)p_3)。
  \newcommand{\barypoint}[4]{%
    \coordinate (#1) at ({6*(#3+0.5*#4)},{6*0.8660254*#4});}
  \coordinate (Wone) at (0,0);
  \coordinate (Wtwo) at (6,0);
  \coordinate (Wthree) at (3,{3*sqrt(3)});
  \barypoint{T1}{.78}{.12}{.10}
  \barypoint{T2}{.10}{.78}{.12}
  \barypoint{T3}{.12}{.13}{.75}
  \barypoint{Doc}{.309}{.4195}{.2715}

  \fill[SLTeal,opacity=.085] (T1)--(T2)--(T3)--cycle;
  \draw[slfig-FIG-P525-01-outer] (Wone)--(Wtwo)--(Wthree)--cycle;
  \node[below left] at (Wone) {$w_1$};
  \node[below right] at (Wtwo) {$w_2$};
  \node[above] at (Wthree) {$w_3$};
  \draw[slfig-FIG-P525-01-hull] (T1)--(T2)--(T3)--cycle;

  \foreach \p/\lab/\pos in {T1/$\phi_{:1}$/left,T2/$\phi_{:2}$/right,T3/$\phi_{:3}$/above right}{
    \fill[SLBlue] (\p) circle (2.6pt) node[\pos,text=SLBlue] {\lab};
  }
  \draw[slfig-FIG-P525-01-spoke] (Doc)--(T1);
  \draw[slfig-FIG-P525-01-spoke] (Doc)--(T2);
  \draw[slfig-FIG-P525-01-spoke] (Doc)--(T3);
  \node[slfig-FIG-P525-01-weight] at (2.05,1.28) {$\theta_{1j}$};
  \node[slfig-FIG-P525-01-weight] at (3.65,.95) {$\theta_{2j}$};
  \node[slfig-FIG-P525-01-weight] at (2.55,2.05) {$\theta_{3j}$};
  \node[diamond,draw=SLGold,fill=SLGold!18,minimum size=6mm,
    inner sep=0pt,line width=.9pt] (dj) at (Doc) {};
  \node[slfig-FIG-P525-01-formula,
    minimum width=58mm,minimum height=8mm] at (8.0,2.55)
    {$P(w\mid d_j)=\sum_{k=1}^{3}\theta_{kj}\phi_{:k}$\\
     $\theta_{kj}\ge0,\ \sum_k\theta_{kj}=1$};
  \node[slfig-FIG-P525-01-legend] at (6.95,1.25)
    {圆：主题点\quad 菱形：文档分布 $P(w\mid d_j)$};
\end{tikzpicture}
\caption{三个单词维度中的主题单纯形：文档分布是主题点按$\theta_{:j}$形成的凸组合}\label{fig:V4-C06-simplex}
\end{figure}

图\ref{fig:V4-C06-simplex}把外层三角形表示为单词单纯形，把蓝色三角形表示为主题凸包。文档点的位置由$\theta_{:j}$决定；不同参数可以给出相同文档点，说明潜在表示不必唯一。
\end{document}
